\documentclass[11pt]{article}

\usepackage[utf8]{inputenc}
\usepackage[T1]{fontenc}
\usepackage{amsmath, amssymb}
\usepackage{graphicx}
\usepackage{booktabs}
\usepackage{hyperref}
\usepackage{geometry}
\usepackage{authblk}
\usepackage{caption}
\usepackage{cite}
\usepackage{setspace}
\usepackage{microtype}
\usepackage{array}
\usepackage{tabularx}
\usepackage{float}

\geometry{
    left=2.5cm,
    right=2.5cm,
    top=2.5cm,
    bottom=2.5cm
}

\title{\textbf{The Evolution of Scientific Claims from Preprint to Publication in AI and Machine Learning: A Large-Scale Study (2015--2024)}}

\author{Moses Boudourides}
\affil{School of Professional Studies, Northwestern University \\ \texttt{Moses.Boudourides@northwestern.edu}}
\date{\today}

\begin{document}

\maketitle

\begin{abstract}
Preprints have become the primary vehicle for early dissemination of research in Artificial Intelligence (AI) and Machine Learning (ML), fundamentally altering the speed and structure of scientific communication. However, the epistemic reliability of preprints relative to their subsequent peer-reviewed publications remains a subject of ongoing debate. In this study, we present a large-scale, automated analysis tracking the evolution of scientific claims across 51,921 AI and ML preprint-publication pairs posted on arXiv, bioRxiv, and medRxiv between 2015 and 2024. Utilizing GPT-4o mini, we extracted and compared claims from abstracts, yielding 217,009 distinct claim-level comparisons. Our findings demonstrate that the majority of claims (59.2\%) remain unchanged through the peer review process. When substantive revisions do occur, claims are significantly more likely to be weakened (14.4\%) than strengthened (11.0\%), reflecting a systematic shift toward more cautious language. At the paper level, 25.3\% of manuscripts were predominantly weakened, whereas only 5.6\% were predominantly strengthened, representing a 4.5:1 ratio of cautious to confident shifts. Furthermore, claim weakening was more pronounced in journal articles compared to conference papers, and manuscripts undergoing longer review periods exhibited a higher relative likelihood of eventual claim strengthening. These results provide robust, large-scale empirical evidence that AI and ML preprints serve as reliable early representations of published research, with the peer review process functioning primarily as a moderating force that constrains overstatement rather than overturning central claims.
\end{abstract}

\section{Introduction}

The landscape of scientific communication has undergone a profound transformation over the past decade, driven largely by the widespread adoption of preprint servers. In the fields of Artificial Intelligence (AI) and Machine Learning (ML), this shift has been particularly dramatic. Platforms such as arXiv have become the de facto standard for the rapid dissemination of novel methodologies, algorithms, and empirical findings, often receiving thousands of submissions per month \cite{wang2020}. This acceleration of the research cycle has facilitated unprecedented rates of innovation and collaboration. Concurrently, the application of AI and ML to the biomedical sciences has fueled a similar surge in preprint postings on bioRxiv and medRxiv \cite{fraser2020}.

Despite their ubiquity, preprints occupy a complex and somewhat ambiguous epistemic status within the broader scientific ecosystem. Because they lack formal peer review prior to public dissemination, preprints carry the inherent risk of propagating overconfident assertions, methodological flaws, or unsupported conclusions that would otherwise be caught and corrected by expert reviewers \cite{rawlinson2020}. Conversely, proponents of the preprint model argue that the benefits of immediate open access---allowing the global scientific community to rapidly evaluate, build upon, or refute new work---far outweigh these risks \cite{berg2016}.

At the heart of this debate lies a crucial empirical question: to what extent do the scientific claims made in preprints change during the peer review process, and in what direction do these changes occur? Understanding the trajectory of claim evolution is essential for determining whether preprints can be trusted as reliable early indicators of scientific evidence, and how heavily they should be weighted by researchers, practitioners, and policymakers.

Previous efforts to quantify the differences between preprints and published articles have largely relied on measures of textual similarity \cite{klein2019, carneiro2020}. While these studies have documented that manuscripts undergo substantial textual revision, such metrics struggle to distinguish between stylistic editing, structural reorganization, and substantive alterations to the core scientific claims. Other approaches have tracked the publication outcomes and citation trajectories of preprints \cite{abdill2019, serghiou2018}, but these macroscopic indicators do not illuminate the granular changes occurring within the text itself.

The recent advent of advanced Large Language Models (LLMs) has opened new avenues for automated, semantic analysis of scientific literature at scale. In a pioneering study, Yin and Rust \cite{yin2026} utilized LLMs to track claim changes across 72,644 biomedical preprints. They found that approximately 90\% of primary claims remained unchanged or were only minorly revised. Crucially, when substantive changes did occur, they shifted toward more cautious language with a 2:1 ratio, suggesting that peer review in the biomedical sciences serves primarily to moderate claims.

In this paper, we extend this analytical framework to the domain of AI and ML. This domain presents a distinct and compelling case study due to its unique publication culture. Unlike the biomedical sciences, where journal publication is the exclusive gold standard, AI and ML research is heavily oriented toward highly competitive, rigorously peer-reviewed conferences (e.g., NeurIPS, ICML, ICLR). Furthermore, the pace of research is exceptionally rapid, and the cultural norms often encourage bold, assertive claims regarding state-of-the-art performance.

To investigate the evolution of claims in this dynamic environment, we analyzed 51,921 AI and ML preprint-publication pairs spanning the years 2015 to 2024. By deploying GPT-4o mini to extract, classify, and compare claims from the abstracts of both the preprint and published versions, we provide the first large-scale, claim-level assessment of preprint reliability in AI and ML. Our analysis addresses not only the frequency and direction of claim changes but also how these patterns vary across preprint sources, publication venues, and over time.

\section{Data and Methods}

\subsection{Corpus Assembly}

The foundation of our analysis is a comprehensive dataset constructed using the Dimensions database, accessed via Google BigQuery (Analytics Hub) and the Dimensions API \cite{dimensions2018}. The Dimensions database is uniquely suited for this task as it maintains extensive linkage metadata connecting preprint records to their subsequently published, peer-reviewed versions.

We initiated our data collection by querying the Dimensions BigQuery dataset for all preprint records posted between January 2015 and December 2024 on arXiv, bioRxiv, and medRxiv that possessed a non-null \texttt{resulting\_publication\_doi} field. This initial query returned 1,085,728 linked preprint-publication pairs. To isolate research within the AI and ML domain, we applied a filtering strategy based on Dimensions' subject classifications (Fields of Research) and a curated list of relevant keywords, which narrowed the dataset to 61,565 candidate pairs.

To ensure the integrity of the preprint-publication links, we performed a validation step using fuzzy string matching. We computed the token sort ratio between the title of the preprint and the title of the published article using the RapidFuzz library. Pairs exhibiting a similarity score below 60 were flagged as likely erroneous matches and removed from the dataset, resulting in the exclusion of 70 pairs (0.11\%). Subsequently, we utilized the Dimensions API to retrieve the full abstracts for both the preprint and published versions of the remaining pairs. We filtered out any pairs where either the preprint or the published abstract was unavailable, empty, or truncated. Finally, we excluded pairs exhibiting a negative time to publication, as these indicate metadata anomalies. The resulting final analysis corpus comprised 51,921 rigorously validated pairs.

\subsection{Claim Extraction and Comparison Pipeline}

To analyze the semantic evolution of the texts, we employed GPT-4o mini \cite{openai2023} to process the paired abstracts. The automated pipeline was designed as a two-stage prompting process to maximize accuracy and structured output.

In the first stage, the model was provided with an abstract and instructed to extract its central scientific claims. The output was required to be a structured list, with each claim accompanied by an assigned certainty label (high, medium, or low) based on the linguistic hedging present in the text.

In the second stage, the model was presented simultaneously with the structured claims extracted from the preprint abstract and those extracted from the published abstract. The model's task was to compare the two sets and classify the relationship of each preprint claim to its published counterpart. The allowable classifications were: \textit{unchanged} (the core semantic meaning and certainty remained identical), \textit{weakened} (the published claim introduced hedging, reduced scope, or lowered certainty), \textit{strengthened} (the published claim removed hedging, expanded scope, or increased certainty), \textit{removed} (the claim appeared in the preprint but was absent from the publication), or \textit{added} (a new claim appeared in the publication that was not present in the preprint). Based on the aggregate of these individual comparisons, the model also assigned a \textit{dominant change type} to the entire pair, representing the plurality of the directional shifts. The full corpus of 51,921 pairs yielded 217,009 claim-level comparisons.

\subsection{Statistical Analysis}

Descriptive statistics were generated to summarize the distribution of change types at both the claim and pair levels. To assess the statistical significance of associations between categorical variables---such as the relationship between dominant change type and preprint source, or change type and publication venue---we employed Chi-square tests of independence.

To formally model the factors influencing the direction of claim revision, we conducted logistic regression analyses using the \texttt{statsmodels} library in Python. We fitted two primary models. The first, a claim-level model, predicted the likelihood of a claim being strengthened versus weakened, conditional on a directional change occurring (n = 55,273). The second, a pair-level model, predicted the likelihood of a manuscript being predominantly strengthened versus predominantly weakened (n = 16,090). Predictor variables included the preprint posting year (centered to aid interpretation), binary indicators for the preprint source (e.g., arXiv vs.\ others), binary indicators for the publication venue type (journal article vs.\ conference paper), and, for the pair-level model, the continuous variable of years from preprint posting to publication.

\section{Results}

\subsection{Corpus Characteristics}

The final analysis corpus of 51,921 preprint-publication pairs is heavily dominated by arXiv, which accounts for 81.8\% of the preprints. The biomedical preprint servers bioRxiv and medRxiv account for 13.6\% and 4.6\%, respectively, reflecting the growing intersection of AI/ML methodologies with the life and health sciences. The mean time from preprint posting to publication across the corpus is 0.63 years. The publication venues are diverse, with journal articles constituting the majority (71.1\%), followed by conference papers (21.9\%) and book chapters (6.0\%). The automated extraction pipeline yielded an average of 4.18 claims per pair, resulting in 217,009 total claim comparisons.

\subsection{The Stability and Moderation of Claims}

Our primary finding is that the scientific claims made in AI and ML preprints are remarkably stable through the peer review process. As shown in Figure \ref{fig:1a}, at the claim level, 59.2\% of all extracted claims remained semantically unchanged between the preprint and the final published version.

\begin{figure}[H]
    \centering
    \includegraphics[width=0.55\textwidth]{fig_1a_claim_change_donut.png}
    \caption{\textbf{Claim-level change from preprint to publication.} Distribution of the 217,009 claim comparisons by change type. The majority of claims (59.2\%) remain semantically unchanged. Among changed claims, weakening (14.4\%) exceeds strengthening (11.0\%), with 10.0\% of preprint claims removed and 5.4\% new claims added in the published version.}
    \label{fig:1a}
\end{figure}

When claims did undergo substantive revision, the direction of that revision was distinctly asymmetric. Claims were significantly more likely to be weakened (14.4\%) than they were to be strengthened (11.0\%), yielding a weakening-to-strengthening ratio of 1.31:1. This systematic shift toward more cautious language is further illustrated in Figure \ref{fig:1b}, which shows the hedging shift at the paper level: 31.1\% of papers moved toward greater caution, while only 20.7\% shifted toward greater confidence, a ratio of 1.5:1.

\begin{figure}[H]
    \centering
    \includegraphics[width=0.55\textwidth]{fig_1b_hedging_shift_donut.png}
    \caption{\textbf{Hedging shift per paper.} For each of the 51,946 pairs, the dominant hedging direction was determined by comparing the number of weakened and strengthened claims. Papers shifting toward greater caution outnumber those shifting toward greater confidence by a ratio of 1.5:1.}
    \label{fig:1b}
\end{figure}

The relationship between the extent of content revision and the direction of hedging shift is examined in Figure \ref{fig:1c}. Papers undergoing major content changes (more than 40\% of claims altered) exhibit the highest proportion of cautious shifts, suggesting that the most substantially revised manuscripts are also those most systematically moderated by peer review.

\begin{figure}[H]
    \centering
    \includegraphics[width=0.7\textwidth]{fig_1c_hedging_by_stratum.png}
    \caption{\textbf{Hedging shift by content-change stratum.} Papers are stratified by the fraction of their claims that underwent substantive change. The proportion shifting toward caution increases markedly with the degree of content revision, while the proportion shifting toward confidence remains comparatively stable.}
    \label{fig:1c}
\end{figure}

The asymmetry between weakening and strengthening becomes even more pronounced when analyzing the data at the paper level. As detailed in Table \ref{tab:pair_level} and illustrated in Figure \ref{fig:1d}, among those papers that exhibited a clear, dominant directional shift, weakening was overwhelmingly prevalent: 25.3\% (13,160 pairs) of manuscripts were classified as predominantly weakened, compared to a mere 5.6\% (2,930 pairs) that were predominantly strengthened, constituting a 4.5:1 ratio. A further 30.4\% (15,764 pairs) remained predominantly unchanged, and 38.1\% showed mixed changes.

\begin{table}[H]
    \centering
    \caption{\textbf{Pair-level dominant change type distribution.} For each of the 51,921 preprint-publication pairs, the dominant change type was determined by the plurality of claim-level classifications.}
    \label{tab:pair_level}
    \begin{tabular}{lrr}
        \toprule
        Dominant Change Type & Pairs & Percentage \\
        \midrule
        Mixed & 19,766 & 38.1\% \\
        Unchanged & 15,764 & 30.4\% \\
        Weakened & 13,160 & 25.3\% \\
        Strengthened & 2,930 & 5.6\% \\
        Removed & 296 & 0.6\% \\
        Added & 4 & 0.0\% \\
        \bottomrule
    \end{tabular}
\end{table}

\begin{figure}[H]
    \centering
    \includegraphics[width=0.85\textwidth]{fig_1d_change_over_time.png}
    \caption{\textbf{Dominant change type over time (pair level), 2015--2024.} The proportions of predominantly weakened, unchanged, strengthened, and mixed papers remain remarkably stable across the decade, with no evidence of a secular trend in either direction.}
    \label{fig:1d}
\end{figure}

\subsection{Variation by Research Field and Claim Type}

The tendency for peer review to moderate rather than amplify claims is consistent across all three preprint sources in the corpus. Figure \ref{fig:1e} compares the hedging shift distributions for arXiv (n = 42,493), bioRxiv (n = 7,068), and medRxiv (n = 2,385). In all three cases, the proportion of papers shifting toward greater caution exceeds the proportion shifting toward greater confidence. The cautious-to-confident ratio is lowest for arXiv (1.48:1), reflecting the greater maturity of arXiv manuscripts at the time of posting, and highest for medRxiv (1.72:1), consistent with the more extensive revision cycles characteristic of clinical research journals.

\begin{figure}[H]
    \centering
    \includegraphics[width=0.75\textwidth]{fig_1e_hedging_by_field.png}
    \caption{\textbf{Hedging shift by preprint source.} For each of the three preprint servers, grouped bars show the percentage of papers shifting toward greater caution (red), exhibiting no hedging shift (grey), or shifting toward greater confidence (green). The cautious-to-confident ratio, annotated above each group, exceeds 1.0 for all three sources, confirming that the moderating effect of peer review is universal. The ratio is lowest for arXiv (1.48:1) and highest for medRxiv (1.72:1).}
    \label{fig:1e}
\end{figure}

The distribution of claim types also shifts between preprint and publication. Figure \ref{fig:1f} presents an alluvial diagram comparing the composition of claim types in preprint abstracts versus their published counterparts. Result claims constitute the largest category in both versions, growing from 79,675 to 83,704, while contribution and conclusion claims also increase modestly. The overall growth in claim count from preprint to publication reflects the addition of new claims in the published versions.

\begin{figure}[H]
    \centering
    \includegraphics[width=0.85\textwidth]{fig_1f_claim_type_alluvial.png}
    \caption{\textbf{Claim-type distribution: preprint versus published.} The alluvial diagram shows the volume of each claim type (result, contribution, conclusion, limitation) in the preprint (left) and published (right) versions. Result claims are the dominant type in both versions, with modest growth across all categories from preprint to publication.}
    \label{fig:1f}
\end{figure}

The weakening rate also varies by claim type, as shown in Figure \ref{fig:2b}. Limitation claims exhibit the highest rate of weakening (17.4\%), followed by result claims (13.1\%). This pattern suggests that peer reviewers are particularly attentive to the scope of stated limitations and the certainty of empirical results.

\begin{figure}[H]
    \centering
    \includegraphics[width=0.75\textwidth]{fig_2b_revision_by_claim_type.png}
    \caption{\textbf{Weakened versus strengthened rate by claim type.} For each major claim type, the percentage of claims that were weakened (red) and strengthened (green) are shown. Limitation claims are weakened most frequently, while all claim types exhibit higher weakening than strengthening rates.}
    \label{fig:2b}
\end{figure}

\subsection{Variations by Source and Venue}

The patterns of claim revision exhibit significant variation depending on both the origin of the preprint and the destination of the publication. Figure \ref{fig:2a} presents a heatmap of dominant change type by preprint source. Preprints originating from arXiv demonstrated the highest rate of unchanged pairs (34.1\%), whereas bioRxiv and medRxiv preprints were substantially more likely to be predominantly weakened (31.0\% and 32.8\%, respectively). These differences are highly statistically significant ($\chi^2 = 4702.57$, df = 8, p < 0.001) and likely reflect the greater maturity of arXiv manuscripts at the time of posting, as well as the more extensive revision cycles characteristic of biomedical journal review.

\begin{figure}[H]
    \centering
    \includegraphics[width=0.75\textwidth]{fig_2a_source_heatmap.png}
    \caption{\textbf{Dominant change type by preprint source.} The heatmap displays the percentage of pairs in each dominant change category, stratified by preprint source. arXiv preprints are more likely to remain unchanged, while bioRxiv and medRxiv preprints are more likely to be predominantly weakened.}
    \label{fig:2a}
\end{figure}

The type of publication venue also exerts a strong influence on the trajectory of claims. As shown in Figure \ref{fig:2e}, the ratio of predominantly strengthened to predominantly weakened papers is below parity for all venue types---0.24 for journal articles, 0.17 for conference papers, and 0.17 for book chapters---confirming that the moderating effect of peer review is universal regardless of venue. However, the magnitude of weakening is greater for journal articles, which exhibit a higher rate of claim weakening (15.8\%) compared to conference papers (11.0\%), while conference papers show a markedly higher rate of unchanged claims (70.7\% versus 55.4\% for journals). These differences ($\chi^2 = 5978.22$, df = 16, p < 0.001) align with the differing structural realities of these review processes.

\begin{figure}[H]
    \centering
    \includegraphics[width=0.65\textwidth]{fig_2e_sw_ratio_by_venue.png}
    \caption{\textbf{Strengthened:Weakened ratio by publication venue type.} All venue types fall below the parity line (dashed), confirming that weakening dominates regardless of where a paper is published. The ratio is lowest for conference papers and book chapters.}
    \label{fig:2e}
\end{figure}

Figure \ref{fig:2f} further disaggregates these patterns by combining source and venue. arXiv papers published in journals show the highest absolute weakening rate among the combinations examined, while arXiv papers published in conference proceedings show the highest rate of claim strengthening (approximately 15\%), consistent with the more assertive norms of the AI conference culture.

\begin{figure}[H]
    \centering
    \includegraphics[width=0.85\textwidth]{fig_2f_change_by_source_venue.png}
    \caption{\textbf{Dominant change type by preprint source and publication venue.} Stacked bars show the proportion of each dominant change type for four major source-venue combinations. arXiv conference papers exhibit the highest strengthening rate, while bioRxiv and medRxiv journal papers show the highest weakening rates.}
    \label{fig:2f}
\end{figure}

\subsection{Temporal Dynamics and Regression Modeling}

Analyzing the data longitudinally reveals that the tendency for peer review to weaken rather than strengthen claims is a highly stable phenomenon. Figure \ref{fig:2c} shows the claim-level change rates from 2015 to 2024. The weakening rate consistently exceeds the strengthening rate, with no evidence of a secular trend in either direction. This stability persists despite the profound methodological shifts in AI over this decade, including the rise of deep learning and transformer architectures.

\begin{figure}[H]
    \centering
    \includegraphics[width=0.85\textwidth]{fig_2c_change_rates_over_time.png}
    \caption{\textbf{Claim change rates over time, 2015--2024.} The weakening rate (red) consistently exceeds the strengthening rate (green) across all years, with the unchanged rate (blue) remaining the dominant category. No secular trend is apparent in any of the three series.}
    \label{fig:2c}
\end{figure}

The relationship between time to publication and the direction of claim revision is examined in Figure \ref{fig:2d}. The weakening rate increases monotonically from 23.5\% for papers published within six months of preprint posting to 28.3\% for papers taking four to eight years. Simultaneously, the strengthening rate rises from 4.7\% to 7.5\% over the same range, though it remains consistently below the weakening rate. The positive trend in weakening rate with time is statistically significant ($R^2 = 0.80$, p = 0.107 for the binned trend line), suggesting that papers subject to extended review are more likely to emerge with moderated claims.

\begin{figure}[H]
    \centering
    \includegraphics[width=0.85\textwidth]{fig_2d_weakening_vs_years.png}
    \caption{\textbf{Weakening and strengthening rates versus time to publication.} Both rates increase with time to publication, but the weakening rate (red) consistently dominates. The dashed line shows the linear trend for weakening rate across publication-time bins.}
    \label{fig:2d}
\end{figure}

To formally isolate the factors associated with claim moderation, we estimated logistic regression models. At the claim level (modeling the 55,273 claims with a directional change), the baseline intercept (OR = 0.724, p < 0.001) formally confirms the overall systemic bias toward weakening. Claims originating from arXiv were found to be slightly more likely to be strengthened relative to weakened when compared to the baseline (OR = 1.094, p < 0.001).

The pair-level model (n = 16,090), which predicts the likelihood of a paper being predominantly strengthened versus predominantly weakened, yielded particularly revealing results. The intercept (OR = 0.174, p < 0.001) underscores the massive 4.5:1 asymmetry at the manuscript level. Controlling for other factors, publication in a journal was associated with a higher relative likelihood of strengthening (OR = 1.261, p = 0.014). Most notably, a longer duration between preprint posting and publication was significantly associated with a higher relative likelihood of the paper emerging with strengthened claims (OR = 1.090 per year, p = 0.001). The full pair-level regression results are detailed in Table \ref{tab:regression_pair}.

\begin{table}[H]
    \centering
    \caption{\textbf{Logistic regression results: pair-level strengthened versus weakened} (n = 16,090). The outcome variable is a binary indicator of predominantly strengthened (1) versus predominantly weakened (0). Odds ratios greater than 1 indicate a higher relative likelihood of strengthening.}
    \label{tab:regression_pair}
    \begin{tabular}{lrrl}
        \toprule
        Predictor & OR & p-value & \\
        \midrule
        Intercept & 0.174 & $<$0.001 & *** \\
        Preprint year (centered) & 1.018 & 0.064 & \\
        arXiv (vs.\ other) & 0.965 & 0.475 & \\
        Journal article (vs.\ other) & 1.261 & 0.014 & * \\
        Conference paper (vs.\ other) & 0.966 & 0.744 & \\
        Years to publication & 1.090 & 0.001 & ** \\
        \bottomrule
    \end{tabular}
\end{table}

\subsection{Journal Prestige and Open-Access Effects}

To examine whether the prestige of the publishing journal or the open-access status of the final publication is associated with the direction of claim revision, we conducted additional analyses using citation counts as a proxy for journal prestige. Figure \ref{fig:3a} plots the weakening rate against the median citation count for 244 journals with at least 20 pairs in the corpus. While there is a statistically significant positive association ($R^2 = 0.03$, p = 0.003), the effect size is modest, suggesting that journal prestige exerts only a marginal influence on the extent of claim moderation.

\begin{figure}[H]
    \centering
    \includegraphics[width=0.75\textwidth]{fig_3a_journal_prestige_scatter.png}
    \caption{\textbf{Journal prestige (median citations) versus weakening rate.} Each point represents a journal with at least 20 pairs in the corpus; bubble size is proportional to the number of pairs. The dashed line shows the linear regression fit. A weak but statistically significant positive association is present, indicating that higher-prestige journals are marginally more likely to weaken claims.}
    \label{fig:3a}
\end{figure}

Figure \ref{fig:3b} stratifies the dominant change type by citation quartile. The weakening rate is remarkably flat across quartiles (approximately 30\% in each), confirming that the moderating function of peer review is not substantially amplified or attenuated by the prestige tier of the publishing venue.

\begin{figure}[H]
    \centering
    \includegraphics[width=0.75\textwidth]{fig_3b_citation_quartile.png}
    \caption{\textbf{Dominant change type by citation quartile.} Papers are divided into four equal quartiles based on the citation count of the published article. The distribution of dominant change types is nearly identical across all four quartiles, indicating that citation impact does not predict the direction of claim revision.}
    \label{fig:3b}
\end{figure}

Finally, Figure \ref{fig:3c} compares the dominant change type for open-access and subscription-access publications. The weakening rates for open-access (25.4\%) and subscription-access (25.1\%) papers are virtually identical, indicating that the access model of the publishing venue has no discernible effect on the direction of claim revision.

\begin{figure}[H]
    \centering
    \includegraphics[width=0.65\textwidth]{fig_3c_open_access.png}
    \caption{\textbf{Dominant change type by open-access status.} Open-access and subscription-access publications exhibit nearly identical distributions of dominant change types, with weakening rates of 25.4\% and 25.1\%, respectively. The access model of the publishing venue does not influence the direction of claim revision.}
    \label{fig:3c}
\end{figure}

\section{Discussion}

The results of this large-scale analysis provide compelling empirical evidence regarding the function of peer review in the fields of Artificial Intelligence and Machine Learning. The central finding---that nearly 60\% of scientific claims remain unchanged from their preprint form to their final published version---offers strong validation for the epistemic reliability of preprints. Researchers, practitioners, and journalists who rely on arXiv, bioRxiv, or medRxiv as early indicators of the state of the art can do so with the confidence that the core empirical claims they encounter are unlikely to be entirely overturned or reversed during formal peer review. This high degree of stability is consistent with recent findings in the biomedical sciences \cite{yin2026}, suggesting that the reliability of preprints is a robust, cross-disciplinary characteristic of the modern scientific ecosystem.

However, when peer review does induce substantive changes, it acts systematically as a moderating force. The pronounced asymmetry we observed---with claims being weakened significantly more often than they are strengthened, culminating in a 4.5:1 ratio at the paper level---highlights the specific value that peer review adds. It functions not primarily to validate or amplify discoveries, but to constrain overstatement, enforce appropriate hedging, and ensure that conclusions do not outpace the underlying evidence.

This moderating effect is likely driven by several intersecting mechanisms. Reviewers routinely identify methodological limitations or edge cases that authors may have overlooked or minimized, prompting revisions that narrow the scope of the claims. Furthermore, the stylistic norms of peer-reviewed venues often demand a more measured, objective tone than the assertive, sometimes promotional language that can characterize preprints, particularly in the highly competitive AI/ML conference landscape. The fact that the 4.5:1 paper-level asymmetry we observe in AI/ML is substantially larger than the 2:1 ratio observed by Yin and Rust \cite{yin2026} in biomedicine suggests that AI/ML preprints may begin from a baseline of greater assertiveness, requiring more aggressive moderation by reviewers.

The structural differences between publication venues further illuminate the mechanics of peer review. The finding that journal articles exhibit higher rates of claim weakening than conference papers is a direct reflection of their differing review paradigms. Journal review is typically an iterative process involving multiple rounds of detailed feedback and revision, providing ample opportunity for claims to be scrutinized and dialed back. In contrast, the rapid, often binary nature of AI conference review---where papers are accepted or rejected with limited scope for major revision---results in a higher proportion of claims passing through unchanged.

Perhaps the most nuanced finding of this study is the relationship between the duration of the review process and the direction of claim revision. Our regression analysis demonstrates that longer times to publication are associated with a higher relative likelihood of claim strengthening. This initially counterintuitive result suggests that extended review periods are not merely indicative of administrative delay or endless rounds of critical pushback. Rather, longer review times likely afford authors the opportunity to conduct additional experiments, gather more data, or perform deeper analyses requested by reviewers. When these extended efforts are successful, they can validate the initial hypotheses more robustly, allowing the authors to ultimately strengthen their claims in the final publication. This dynamic implies that the rapid-turnaround conference model, while highly efficient for dissemination, may systematically undervalue the epistemic benefits---including the potential for claim fortification---that arise from thorough, extended peer review.

The absence of any meaningful association between journal prestige (as proxied by citation counts) and the direction of claim revision is also theoretically significant. It suggests that the moderating function of peer review is not a privilege of elite venues but a consistent feature of the scholarly communication system at large. Similarly, the finding that open-access and subscription-access publications exhibit identical revision patterns indicates that the access model of the venue does not introduce any systematic bias into the review process.

\section{Limitations}

While this study leverages a massive dataset to provide unprecedented visibility into claim evolution, several limitations must be acknowledged. First, our analysis is strictly confined to the abstracts of the manuscripts. While abstracts are designed to distill the most critical claims of a paper, they inherently omit the granular details of methodology, statistical analysis, and nuanced argumentation found in the full text. Therefore, our results represent a lower bound on the total volume of revision; papers that appear unchanged in their abstracts may have undergone substantial methodological refinement in the body text.

Second, the study relies entirely on an automated LLM pipeline for claim extraction and classification. Although models like GPT-4o mini demonstrate high proficiency in natural language understanding tasks, automated classification of complex, highly technical scientific prose is susceptible to error and misinterpretation. A human-in-the-loop validation study is currently underway, in which two independent domain experts will review a stratified random sample of 200 pairs to assess the accuracy of the claim extraction and the correctness of the change classifications. Inter-rater reliability will be assessed using Cohen's kappa.

Third, the analysis is subject to survivorship bias. We examine only those preprints that successfully navigated the peer review process and achieved publication. Preprints that were rejected, fundamentally flawed, or abandoned are absent from this corpus. Consequently, our findings describe the evolutionary trajectory of viable research, and our estimates of claim stability cannot be extrapolated to the entire universe of posted preprints.

Fourth, the classification of publication venue types relies on metadata from the Dimensions database, which may contain inaccuracies, particularly for venues that publish both proceedings and journal versions of the same work.

Fifth, the study covers the period 2015 to 2024, which largely predates the widespread use of LLMs as writing assistants in scientific research. The growing use of tools such as ChatGPT in manuscript preparation may alter the relationship between preprint and published abstracts in ways that are not captured by our analysis, and future work should examine whether the patterns we observe change as LLM-assisted writing becomes more prevalent.

\section{Conclusion}

This study provides the first large-scale, quantitative assessment of how scientific claims in Artificial Intelligence and Machine Learning evolve from preprint to peer-reviewed publication. Our analysis of nearly 52,000 manuscript pairs reveals a dual reality: preprints are highly reliable early indicators of published research, with the majority of claims remaining stable, yet peer review performs a vital, systematic function in moderating overstatement. When claims are revised, they are overwhelmingly shifted toward greater caution and precision, and this moderating effect is consistent across research fields, preprint sources, publication venues, journal prestige tiers, and the entire decade from 2015 to 2024. As the volume of AI and ML research continues to accelerate, and as preprints increasingly drive both academic discourse and public policy, understanding this dynamic is essential. Preprints can be trusted to convey the core direction of research, but readers must anticipate that the final, peer-reviewed record will likely present a more measured, carefully bounded version of those initial discoveries.

\bibliographystyle{unsrt}
\bibliography{references}

\end{document}
